\clearpage
\section{Projektowanie}

\subsection{Podział odpowiedzialności}

Program składa się z interfejsu konsolowego, klasy reprezentującej
drzewo oraz modułu operacji plikowych. Tabela~\ref{tab:moduly}
pokazuje, za co odpowiada każda część. Deklaracje struktur i metod
umieszczono w plikach nagłówkowych. Ich szczegółowy wykaz zawiera
również dokumentacja wygenerowana przez Doxygen~\cite{doxygen-bst}.

\begin{tabela}{Podział programu na moduły}
              {Podział programu na moduły}
              {\label{tab:moduly}}
\small
\begin{tabularx}{\linewidth}{|L{4.2cm}|Y|}
\hline
\textbf{Pliki} & \textbf{Odpowiedzialność}\\ \hline
\kod{main.cpp} & Wyświetla menu, sprawdza dane wejściowe i przekazuje
                 operacje do \kod{BST} lub \kod{BSTFile}.\\ \hline
\kod{BST.h}, \kod{BST.cpp} & Definiują węzeł \kod{Node} i klasę
  \kod{BST}. Zarządzają węzłami, wyszukiwaniem, przejściami,
  rysunkiem drzewa i jego zapisem tekstowym.\\ \hline
\kod{BSTFile.h}, \kod{BSTFile.cpp} & Zawierają metody statyczne
  do zapisu i odczytu binarnego oraz do importu tekstowego.
  Korzystają z metod publicznych klasy \kod{BST}.\\ \hline
\kod{CMakeLists.txt} & Definiuje budowanie programu
  \kod{Drzewo_BST} z trzech plików \kod{.cpp}.\\ \hline
\end{tabularx}
\end{tabela}

\subsection{Warunki poprawności drzewa}

Wektor \kod{tree} zawiera wszystkie węzły należące do drzewa.
W drzewie niepustym \kod{rootIndex} wskazuje prawidłową pozycję
w wektorze. Każdy pozostały węzeł jest osiągalny z korzenia przez
jedno połączenie od rodzica. Pola dzieci przyjmują wartość $-1$
albo indeks istniejącego węzła. Przez całą ścieżkę obowiązuje
porządek BST opisany w rozdziale~\ref{sec:analiza}.

Indeks jest szczegółem reprezentacji i może zmienić się po
usunięciu pozycji z wektora. Z tego powodu metody publiczne
przyjmują klucze typu \kod{int}, a \kod{findPath()} i
\kod{preorderValues()} zwracają wartości węzłów. Program główny
może więc używać tych metod bez dostępu do wektora i indeksów.

\Needspace{11\baselineskip}
\subsection{Menu konsolowe}
\label{sec:menu}

Opcje menu zestawiono w tabeli~\ref{tab:menu}. W trybach dodawania
i usuwania można podać wiele wartości rozdzielonych białymi znakami.
Wpisanie małej litery \kod{q} przywraca menu. Wyszukiwanie ścieżki
przyjmuje jedną wartość i po wyświetleniu wyniku wraca do menu.

\begin{tabela}{Opcje menu programu}{Opcje menu programu}{\label{tab:menu}}
\small
\begin{tabularx}{\linewidth}{|c|Y|L{5.1cm}|}
\hline
\textbf{Opcja} & \textbf{Działanie} & \textbf{Metoda}\\ \hline
1 & Dodawanie elementów & \kod{BST::addNode()}\\ \hline
2 & Usuwanie elementów & \kod{BST::removeNode()}\\ \hline
3 & Usuwanie całego drzewa & \kod{BST::removeTree()}\\ \hline
4, 5, 6 & Przejścia preorder, inorder i postorder.
Każda opcja wypisuje również rysunek drzewa. &
\kod{preorder()}, \kod{inorder()}, \kod{postorder()}, \kod{printTree()}\\ \hline
7 & Wyszukanie ścieżki od korzenia & \kod{BST::findPath()}\\ \hline
8 & Zapis do \kod{drzewo.txt} & \kod{BST::saveToTextFile()}\\ \hline
9 & Zapis do \kod{drzewo.bin} & \kod{BSTFile::saveBinary()}\\ \hline
10 & Odtworzenie drzewa z \kod{drzewo.bin} & \kod{BSTFile::loadBinary()}\\ \hline
11 & Dodanie liczb ze wskazanego pliku tekstowego & \kod{BSTFile::loadText()}\\ \hline
12 & Zakończenie programu & Powrót z \kod{main()}\\ \hline
\end{tabularx}
\end{tabela}

Numery opcji i wartości są odczytywane jako pojedyncze tokeny.
Program sprawdza, czy cały token można odczytać jako liczbę typu
\kod{int}. Wpisy takie jak \kod{3x}, \kod{3.5} i liczby spoza
zakresu typu są odrzucane. Liczba całkowita spoza zakresu opcji
1--12 powoduje ponowne wyświetlenie menu z komunikatem o błędnej
opcji. Zakończenie strumienia wejściowego kończy pracę programu.

\Needspace{9\baselineskip}
\subsection{Dane zapisywane w plikach}
\label{sec:formaty}

Zapis tekstowy tworzy plik \kod{drzewo.txt} z jedną liczbą w każdym
wierszu. Zapis binarny tworzy plik \kod{drzewo.bin}, który zaczyna się
od licznika typu \kod{std::uint64_t}, po którym występują kolejne
wartości typu \kod{int}. W obu przypadkach wartości są pobierane
w kolejności preorder. Plik binarny zapisany przez program zawiera
więc $8 + n \cdot \texttt{sizeof(int)}$ bajtów dla $n$ węzłów.
Program zapisuje typy w ich reprezentacji maszynowej. Przy wymianie
plików między platformami wymagane są zgodne rozmiary i kolejność
bajtów zapisywanych typów.

Wstawianie wartości w kolejności preorder do pustego drzewa
odtwarza jego kształt. Pierwsza wartość staje się korzeniem.
Przed wartościami każdego poddrzewa występuje jego własny korzeń,
więc ta sama zasada działa na kolejnych poziomach. Węzły w wektorze
mogą po odczycie mieć inne indeksy, ponieważ wynikowy wektor jest
budowany w kolejności wstawiania.

Odczyt binarny zastępuje bieżące drzewo wynikiem po sprawdzeniu
całego pliku. Import tekstowy dopisuje unikatowe liczby do
bieżącego drzewa. W opcji~11 należy podać ścieżkę w jednym
wierszu. Ścieżka może zawierać spacje i jest podawana bez
cudzysłowów. Ścieżki względne są rozwiązywane względem katalogu
roboczego programu. Jeśli celem jest odtworzenie drzewa z pliku
\kod{drzewo.txt}, należy najpierw je opróżnić opcją~3.
